\chapter{测度空间}
\label{chap:measure-spaces}

Riemann 积分把区间分成有限多个小块，再让分割变细。这个构造适合边界较
规则的集合，却不能稳定处理可数次并、可数次交和逐点极限。测度论从集合
出发：先规定哪些集合可以测量，再要求长度、面积或质量对可数个互不相交
部分相加。可数性并非形式上的强化；极限过程天然产生可数个集合。

本章先在抽象集合上建立这套语言。Lebesgue 测度将在后面由外测度构造，
这里不预设集合来自欧氏空间。

\section{可测集合族}

设 $X$ 是非空集合，$f_n,f:X\to\R$ 为有限实值函数。若只要求集合族
对有限并和补集封闭，便不足以处理例如
\[
 \{x\in X:f_n(x)\to f(x)\}
 =\bigcap_{k=1}^{\infty}\bigcup_{N=1}^{\infty}
   \bigcap_{n\ge N}\{x\in X:|f_n(x)-f(x)|<1/k\}
\]
这样的集合。有限复值函数也满足此式，扩充值函数则不能直接写可能出现的
$\infty-\infty$。可数运算确实增加了要求：在 $\N$ 上，所有有限集和有限集
的补集组成的集合族对有限并和补集封闭，却不包含偶数集；后者是可数个
单点集的并。为了处理上述极限集合，需要更强的封闭性。

\begin{definition}
集合族 $\mathcal A\subseteq\mathcal P(X)$ 称为 $X$ 上的
\term{$\sigma$-代数}，如果
\begin{enumerate}[label=\textup{(\roman*)}]
 \item $X\in\mathcal A$；
 \item $E\in\mathcal A$ 时 $E^c=X\setminus E\in\mathcal A$；
 \item $E_n\in\mathcal A$（$n\in\Npos$）时
       $\bigcup_{n=1}^{\infty}E_n\in\mathcal A$。
\end{enumerate}
二元组 $(X,\mathcal A)$ 称为\term{可测空间}，$\mathcal A$ 的成员称为
\term{可测集}。
\end{definition}

由 De Morgan 公式，$\sigma$-代数也对可数交封闭；取常值集合列还可得到
$\varnothing\in\mathcal A$ 以及有限并、有限交和差的封闭性。最小的
$\sigma$-代数是 $\{\varnothing,X\}$，最大的则是 $\mathcal P(X)$。

任意多个 $X$ 上的 $\sigma$-代数的交仍为 $\sigma$-代数，空族的交相对于
$\mathcal P(X)$ 取，约定为 $\mathcal P(X)$。因此，对任意
$\mathcal E\subseteq\mathcal P(X)$，包含 $\mathcal E$ 的最小
$\sigma$-代数
\[
 \sigma(\mathcal E)
 :=\bigcap\{\mathcal A\in\mathcal P(\mathcal P(X)):\mathcal A\text{ 是包含 }\mathcal E
 \text{ 的 }\sigma\text{-代数}\}
\]
定义良好，称为由 $\mathcal E$ \term{生成}的 $\sigma$-代数。这里的交集
描述通常比逐项列出成员更有用；即使 $\mathcal E$ 很简单，生成结果也可能
包含复杂的可数运算。

拓扑空间 $X$ 的开集生成的 $\sigma$-代数记为 $\mathcal B(X)$，称为
\term{Borel $\sigma$-代数}。在 $\R$ 中，有理端点开区间已经足以生成
$\mathcal B(\R)$：每个开集是包含于它的有理端点开区间之并，这个族至多
可数；反过来，这些区间本来就是开集。

所有半直线 $(-\infty,a)$ 也生成同一个集合族。它们都是开集，而
\[
 (-\infty,u]=\bigcap_{n=1}^{\infty}(-\infty,u+1/n),\qquad
 (u,v)=(-\infty,v)\setminus(-\infty,u]
\]
说明它们生成所有开区间，进而生成所有开集。若改用所有有界闭区间，结果
仍相同：闭区间是 Borel 集；每个开区间 $(u,v)$ 又是闭区间
$[u+1/n,v-1/n]$ 的可数并，只取使左端不大于右端的正整数 $n$。
这些论证同时核对了两个方向的包含。

\section{测度与可数可加性}

以下统一使用扩充非负实数 $[0,+\infty]$。约定 $0\cdot(+\infty)=0$ 只在
以后简单函数积分的无歧义语境中使用；测度的可加性不需要这一约定。

\begin{definition}
设 $(X,\mathcal A)$ 为可测空间。映射
$\mu:\mathcal A\to[0,+\infty]$ 称为\term{测度}，如果
\begin{enumerate}[label=\textup{(\roman*)}]
 \item $\mu(\varnothing)=0$；
 \item 对两两不交的可测集列 $(E_n)$，有
 \[
  \mu\!\left(\bigcup_{n=1}^{\infty}E_n\right)
  =\sum_{n=1}^{\infty}\mu(E_n).
 \]
\end{enumerate}
三元组 $(X,\mathcal A,\mu)$ 称为\term{测度空间}。
\end{definition}

若 $\mu(X)<\infty$，称 $\mu$ 为有限测度。若存在 $X_n\in\mathcal A$ 使
$X=\bigcup_nX_n$ 且 $\mu(X_n)<\infty$，称 $\mu$ 为
\term{$\sigma$-有限测度}。第二个条件允许把许多有限测度论证分块使用，
但它不能在未说明时自动加入定理假设。

\begin{example}
在任意集合 $X$ 上，计数测度定义为
\[
 \#E=\begin{cases}|E|,&E\text{ 有限},\\+\infty,&E\text{ 无限}.
 \end{cases}
\]
它定义在 $\mathcal P(X)$ 上。固定 $a\in X$，Dirac 测度
$\delta_a(E)=\mathbf1_E(a)$ 定义在 $\mathcal P(X)$ 上，把全部质量集中在一点。
更一般地，把至多可数集合无重复地写成 $X=\{x_k:k\in I\}$，其中 $I$
是有限指标集或 $\Npos$。给定权重 $w_k\in[0,+\infty]$，公式
\[
 \mu(E)=\sum_{\{k\in I:x_k\in E\}}w_k
\]
给出加权计数测度。这里按属于集合的指标求和，空和为零，不引入
$0\cdot\infty$。对不交集合列，每个指标至多落在一个集合中；非负项之和
等于所有有限子和的上确界，所以先按集合分组或直接对全部指标求和结果相同。
这就验证了可数可加性，并未使用积分的交换定理。
\end{example}

计数测度的可数可加性也可直接核对：不交并有限时逐点计数；若并无限，
要么某一项无限，要么有限部分并的点数没有上界，两种情形下右侧均为
$+\infty$。与长度不同，计数测度对单个点就赋予质量 $1$，可见测度的定义
规定的是加法规则，并不预先指定“大小”必须是哪一种几何量。

\section{测度的基本运算律}

\begin{theorem}[单调性、次可加性与连续性]
\label{thm:measure-basic-continuity}
设 $(X,\mathcal A,\mu)$ 为测度空间。
\begin{enumerate}[label=\textup{(\roman*)}]
 \item 若 $E,F\in\mathcal A$ 且 $E\subseteq F$，则 $\mu(E)\le\mu(F)$；若另有
 $\mu(E)<\infty$，则 $\mu(F\setminus E)=\mu(F)-\mu(E)$。
 \item 对任意可测集列 $(E_n)$，
 \[
  \mu\!\left(\bigcup_nE_n\right)\le\sum_n\mu(E_n).
 \]
 \item 若 $E_n\uparrow E$，即 $E_n\subseteq E_{n+1}$ 且
 $E=\bigcup_nE_n$，则 $\mu(E_n)\uparrow\mu(E)$。
 \item 若 $E_n\downarrow E$ 且某个 $E_{n_0}$ 的测度有限，则
 $\mu(E_n)\downarrow\mu(E)$。
\end{enumerate}
\end{theorem}

\begin{proof}
由 $F=E\mathbin{\dot\cup}(F\setminus E)$ 得
$\mu(F)=\mu(E)+\mu(F\setminus E)$，于是得到 (i)；有限性正是允许从等式
两端减去 $\mu(E)$ 的条件。

对 (ii)，令
\[
 F_1=E_1,\qquad
 F_n=E_n\setminus\bigcup_{k<n}E_k\quad(n\ge2).
\]
则 $F_n$ 两两不交、$F_n\subseteq E_n$，且 $\bigcup_nF_n=\bigcup_nE_n$。
可数可加性和单调性给出结论。

若 $E_n\uparrow E$，令 $F_1=E_1$、$F_n=E_n\setminus E_{n-1}$。这些
集合两两不交，且 $E_n=\bigcup_{k\le n}F_k$、
$E=\bigcup_kF_k$，故
\[
 \mu(E_n)=\sum_{k\le n}\mu(F_k)
 \longrightarrow\sum_{k=1}^{\infty}\mu(F_k)=\mu(E).
\]
这证明 (iii)。

对 (iv)，删去有限项后可设 $\mu(E_1)<\infty$。集合
$E_1\setminus E_n$ 递增到 $E_1\setminus E$，由 (iii)
\[
 \mu(E_1)-\mu(E_n)
 \longrightarrow\mu(E_1)-\mu(E).
\]
这里各项有限，移项即得所需结论。
\end{proof}

下降连续性中的有限性不可删除。对计数测度取
$E_n=\{n,n+1,\ldots\}\subset\N$，则 $E_n\downarrow\varnothing$，
但每个 $\mu(E_n)=+\infty$。

\section{完备化}

测度为零的集合在积分中不产生贡献。若一个零测集的子集不在原
$\sigma$-代数中，后续每次修改函数的零测值都要重新检查可测性。完备化
把这些子集一次性补入。

\begin{definition}
若对每个满足 $N\in\mathcal A$、$\mu(N)=0$ 的集合 $N$，其每个子集
$Z\subseteq N$ 都属于 $\mathcal A$，称测度空间 $(X,\mathcal A,\mu)$ \term{完备}。
\end{definition}

\begin{theorem}[测度空间的完备化]
\label{thm:measure-completion}
令 $\mathcal N\subseteq\mathcal P(X)$ 为全部可测零测集的子集所成的集合族，并令
\[
 \overline{\mathcal A}
 =\{E\cup Z:E\in\mathcal A, Z\in\mathcal N\}.
\]
则 $\overline{\mathcal A}$ 是 $\sigma$-代数，公式
\[
 \overline\mu(E\cup Z)=\mu(E)
 \qquad(E\in\mathcal A, Z\in\mathcal N)
\]
定义了一个完备测度。它延拓 $\mu$，并且是包含 $\mathcal A$ 的最小完备
延拓。
\end{theorem}

\begin{proof}
先把集合族改写成对称差形式。若 $A=E\cup Z$，其中
$Z\subseteq N\in\mathcal A$ 且 $\mu(N)=0$，则
$A\mathbin\triangle E\subseteq N$。反过来，若
$A\mathbin\triangle E\subseteq N$，则
\[
 A=(E\setminus N)\cup(A\cap N),
\]
其中 $E\setminus N\in\mathcal A$，而 $A\cap N\in\mathcal N$。因此
\begin{equation}
 A\in\overline{\mathcal A}
 \quad\Longleftrightarrow\quad
 \text{存在 }E\in\mathcal A\text{ 及可测零集 }N,
 \ A\mathbin\triangle E\subseteq N.
 \label{eq:completion-symmetric-difference}
\end{equation}

由这个刻画验证集合运算。若
$A\mathbin\triangle E\subseteq N$，则
$A^c\mathbin\triangle E^c=A\mathbin\triangle E\subseteq N$。若
$A_k\mathbin\triangle E_k\subseteq N_k$，则
\[
 \left(\bigcup_kA_k\right)\mathbin\triangle
 \left(\bigcup_kE_k\right)
 \subseteq\bigcup_kN_k.
\]
右端是可测零集。故 $\overline{\mathcal A}$ 对补集和可数并封闭，因而是
$\sigma$-代数。

还须核对测度值不依赖可测代表。若同一集合 $A$ 分别与
$E,F\in\mathcal A$ 的对称差分别包含于原可测零集，则 $E\mathbin\triangle F$ 也包含于某个
可测零集。于是
$\mu(E\setminus F)=\mu(F\setminus E)=0$，从而
$\mu(E)=\mu(E\cap F)=\mu(F)$。所以定义
$\overline\mu(A)=\mu(E)$ 是良定义的。

设 $A_k\in\overline{\mathcal A}$ 两两不交，并取可测代表 $E_k$，使
$A_k\mathbin\triangle E_k\subseteq N_k$，其中 $N_k\in\mathcal A$ 且
$\mu(N_k)=0$。令
\[
 F_1=E_1,\qquad
 F_k=E_k\setminus\bigcup_{j<k}E_j\quad(k\ge2).
\]
$F_k$ 两两不交且可测。由于 $A_j\cap A_k=\varnothing$，交集
$E_j\cap E_k$ 只能落在代表误差的并中，因而为零测集；所以
$E_k\mathbin\triangle F_k$ 零测，进而
$\mu(F_k)=\mu(E_k)=\overline\mu(A_k)$。同时
$\bigcup_kA_k$ 与 $\bigcup_kF_k$ 的对称差包含于可测零集
$\bigcup_kN_k$，故
\[
 \overline\mu\left(\bigcup_kA_k\right)
 =\mu\left(\bigcup_kF_k\right)
 =\sum_k\mu(F_k)
 =\sum_k\overline\mu(A_k).
\]

若 $Z\subseteq A\in\overline{\mathcal A}$ 且 $\overline\mu(A)=0$，取可测
代表 $E$ 及原可测零集 $N$，使 $A\triangle E\subseteq N$。
由 $\mu(E)=\overline\mu(A)=0$，$A\cup E\subseteq E\cup N$，右端是原可测零集，因此 $Z$ 也属于
$\overline{\mathcal A}$；这证明完备性。最后，任一包含 $\mathcal A$ 且
延拓 $\mu$ 的完备测度空间都必须包含所有可测零集的子集，继而包含每个
$E\cup Z$。故上述完备化在集合域意义下最小。
\end{proof}

\section{集合列的上下极限}

对集合列 $(E_n)_{n\ge1}$ 定义
\[
 \liminf_{n\to\infty}E_n
 =\bigcup_{N=1}^{\infty}\bigcap_{n\ge N}E_n,
 \qquad
 \limsup_{n\to\infty}E_n
 =\bigcap_{N=1}^{\infty}\bigcup_{n\ge N}E_n.
\]
集合 $\liminf E_n$ 由最终属于全部 $E_n$ 的点组成；集合
$\limsup E_n$ 由属于无穷多个 $E_n$ 的点组成。总有
$\liminf E_n\subseteq\limsup E_n$。二者相等时，指标函数列
$\mathbf1_{E_n}$ 逐点收敛。

\begin{proposition}[集合列的测度估计]
\label{prop:set-liminf-limsup-measure}
对任意可测集列，
\[
 \mu(\liminf E_n)\le\liminf_n\mu(E_n).
\]
若存在 $N$ 使 $\mu(\bigcup_{n\ge N}E_n)<\infty$，则
\[
 \limsup_n\mu(E_n)\le\mu(\limsup E_n).
\]
特别地，若 $\sum_n\mu(E_n)<\infty$，则
$\mu(\limsup E_n)=0$。
\end{proposition}

\begin{proof}
令 $F_N=\bigcap_{n\ge N}E_n$。则 $F_N\uparrow\liminf E_n$，而
$\mu(F_N)\le\inf_{n\ge N}\mu(E_n)$；用上升连续性并令 $N\to\infty$
得到第一式。令 $G_N=\bigcup_{n\ge N}E_n$，则
$G_N\downarrow\limsup E_n$。在给定有限性条件下使用下降连续性，且
$\mu(G_N)\ge\sup_{n\ge N}\mu(E_n)$，得到第二式。最后，次可加性给出
$\mu(G_N)\le\sum_{n\ge N}\mu(E_n)\to0$。
\end{proof}

\section*{习题}
\addcontentsline{toc}{section}{习题}

\subsection*{A组　可测集合与测度运算}
\begin{enumerate}[label=\textbf{\thechapter.\arabic*.}]
 \exerciseitem{sigma-algebra-intersections}
 证明任意多个 $\sigma$-代数的交仍为 $\sigma$-代数，并说明它们的并通常
 不是 $\sigma$-代数。
 \exerciseitem{finite-generator-atoms}
 设 $m\in\N$、$\mathcal E=\{E_1,\ldots,E_m\}\subseteq\mathcal P(X)$。
 证明 $\sigma(\mathcal E)$ 由至多 $2^m$ 个非空交块
 $\bigcap_{j=1}^mF_j$ 的任意并组成，其中 $F_j=E_j$ 或 $E_j^c$。
 证明这些交块是原子，即不含该 $\sigma$-代数的非空真子集。
 $m=0$ 时，空交约定为 $X$。
 \exerciseitem{measure-finite-additivity}
 从可数可加性推出有限可加性、单调性和有限并的次可加不等式。
 \exerciseitem{continuity-from-below-proof}
 不引用定理 \ref{thm:measure-basic-continuity}，用两两不交分解证明测度的
 上升连续性。
 \exerciseitem{continuity-from-above-counterexample}
 构造下降集合列说明下降连续性的有限测度条件不能删除。
\end{enumerate}

\subsection*{B组　典型测度与完备化}
\begin{enumerate}[label=\textbf{\thechapter.\arabic*.},resume]
 \exerciseitem{weighted-counting-measure}
 验证加权计数公式定义测度，并判定它何时有限、何时 $\sigma$-有限。
 \exerciseitem{dirac-integral-preview}
 设 $a\in X$，$\delta_a$ 定义在 $\mathcal P(X)$ 上。证明它是完备测度，并说明
 $\delta_a(E_n)\to\delta_a(E)$ 等价于指标函数在 $a$ 点收敛。
 \exerciseitem{completion-equivalent-form}
 证明 $A\in\overline{\mathcal A}$ 当且仅当存在 $E,F\in\mathcal A$ 使
 $E\subseteq A\subseteq F$ 且 $\mu(F\setminus E)=0$。
 \exerciseitem{completion-minimality}
 完成定理 \ref{thm:measure-completion} 的最小性证明。
 \exerciseitem{measure-algebra-metric}
 若 $\mu(X)<\infty$，在可测集按零测对称差所得的等价类上定义
 $d([E],[F])=\mu(E\mathbin\triangle F)$。证明这是度量。
\end{enumerate}

\subsection*{C组　集合列}
\begin{enumerate}[label=\textbf{\thechapter.\arabic*.},resume]
 \exerciseitem{indicator-set-limits}
 证明 $\mathbf1_{E_n}$ 收敛当且仅当
 $\liminf E_n=\limsup E_n$，并写出极限的指标函数。
 \exerciseitem{set-fatou-inequality}
 直接从上升连续性证明
 $\mu(\liminf E_n)\le\liminf\mu(E_n)$。
 \exerciseitem{first-borel-cantelli}
 设 $E_n\in\mathcal A$。证明若 $\sum_n\mu(E_n)<\infty$，则几乎每个点
 只属于有限多个 $E_n$。这里“几乎每个”指除去一个可测零集之后的每个点。
 \optionalexerciseitem{pushforward-measure}
 若 $T:X\to Y$ 且 $\mathcal B$ 是 $Y$ 上的 $\sigma$-代数，假设
 对每个 $B\in\mathcal B$ 都有 $T^{-1}(B)\in\mathcal A$。证明
 $T_\#\mu(B)=\mu(T^{-1}(B))$ 定义 $Y$ 上的测度。
\end{enumerate}
